All 1,499 maintained tests pass in 258.931 seconds. The 66 focused integration
tests pass in 9.620 seconds. They retain the delivered geometric, integral
transport, callback, causal projection and composition checks, and add
actual recognizer/CLI positive replay, cap-safe complete fallback, unchanged
disabled behavior, a source-first index bypass and an exhausted-node index
bypass. The report's historical endpoints are read from its unchanged
artifacts rather than copied into a new search result.

The native source and diagram audit completes in 41.188 seconds. Ten copied
runtime modules and six foundational move/cocycle/codec modules exactly
match their delivered source bytes. The diagram adapter is the reconciled
part: bounded sleep/naive queries use shared covers after the first source
probe, and the recognizer and CLI provide the optional stage.

The two genuine source fixtures attach one or two six-cell descent gadgets
to a finite solid torus, giving seven and thirteen initial tetrahedra. Neither
has an initial native $3\to2$ move. Independent restarted and shared
exhaustions at $U=1$, cover cardinality six, have identical exact
cochain-isomorphism endpoint classes: fourteen and twenty-seven. The native
shared engine retains twenty-eight and fifty-five marked endpoints, versus
495 and 3,152 restarted frames. All 124 retained source endpoint proofs
pass independent replay. The exact class comparisons use full canonical
cochain-isomorphism tuples, not digest-only equality; they are distinct from
the runtime's marked incarnation keys.

At $U=0$ both sources completely exhaust bounded descent. At $U=1$ both
return a separately replayed one-up/two-down witness with strict loss one.
Fresh Regina checks confirm that each original and returned endpoint is a
solid torus. These checks do not replace the source-bound move replay or
prove a universal small-upward theorem.

The first-descent comparison uses five randomized paired rounds and four
arms per source, including identical-baseline A/A controls. Forty measured
calls and eight warmups complete their configured positive query, independent
replay and full serialization. Both arms validate and search the same source
and declared $U=1$ connected-cover family. Source/index preparation and the
actual transitions are charged in both arms.

\begin{center}
\small
\begin{tabular}{lrrrrr}
\toprule
Input & previous ms & current ms & old/new & old A/A & new A/A\\
\midrule
descent 1 & 989.355 & 26.204 & 37.721 & 0.991 & 1.014\\
descent 2 & 3862.141 & 48.123 & 79.691 & 0.997 & 0.995\\
\bottomrule
\end{tabular}
\end{center}


The median paired old/new ratios are 37.721 and 79.691. The first case
begins 273 versus six frames and records 210,418 versus 5,263 callback work
units; the second begins 630 versus six frames and records 837,402 versus
14,520. Each returned proof gives strict loss one. The A/A controls stay
within about two percent of parity. These measurements establish gains for
complete \emph{first-descent source queries}, not full family exhaustion or
complete knot recognition on these sources.

The diagram audit uses thirteen actual diagrams, including unknots,
trefoil and figure-eight controls and hard native-seed survivors. With a
200,000-unit local allowance, all thirteen complete enabled recognition
verdicts agree with the frozen preceding recognizer. Disabled statuses,
methods and evidence agree exactly. The direct diagram adapter certifies
the empty-circle case; all twelve other native queries reach the work cap.
The one saved positive also receives fresh Regina confirmation. No added
nonempty-diagram coverage is claimed for this cohort. Separate focused tests
certify nonempty one-crossing inputs, both zero-move and transported positive
proofs, and binding against a different input knot.

Five paired rounds complete 180 measured native/full-recognition calls and
36 warmups. Native query rows compare the delivered independently restarted
adapter with the maintained shared adapter, both under the same local work
cap. Full recognition rows compare enabling this new stage with omitting it.
Preparation, proof replay, serialization and any local cap/fallback are
included. The four nonempty native timing controls all cap, so their ratios
are bounded-call latency observations rather than completed-family speedups.

\textbf{Bounded native diagram queries.}
\begin{center}
\small
\begin{tabular}{lrrrrr}
\toprule
Input & previous ms & current ms & old/new & old A/A & new A/A\\
\midrule
empty circle & 9.994 & 9.711 & 1.029 & 1.001 & 0.992\\
optimized positive & 788.913 & 191.477 & 4.120 & 0.995 & 1.002\\
genus one miss & 883.092 & 200.690 & 4.516 & 1.017 & 0.993\\
trefoil & 873.501 & 184.554 & 4.733 & 0.998 & 0.966\\
figure eight & 861.689 & 187.937 & 4.569 & 1.005 & 0.990\\
\bottomrule
\end{tabular}
\end{center}

\textbf{Complete recognition, disabled versus enabled Pachner stage.}
\begin{center}
\small
\begin{tabular}{lrrrrr}
\toprule
Input & previous ms & current ms & old/new & old A/A & new A/A\\
\midrule
empty circle & 0.016 & 0.016 & 1.018 & 0.975 & 0.975\\
genus one miss & 0.687 & 228.249 & 0.003 & 1.087 & 1.060\\
trefoil & 0.637 & 211.157 & 0.003 & 0.660 & 0.873\\
figure eight & 0.900 & 221.463 & 0.004 & 1.026 & 1.012\\
\bottomrule
\end{tabular}
\end{center}


The capped native queries have old/new ratios about 4.12--4.73, while the
source-positive empty-circle query is near parity. This does not establish
successful search on the capped diagrams. Enabling the stage adds roughly
0.21--0.23 seconds to the three small nonempty recognition controls whose
complete fallback takes under one millisecond. Their paired disabled/enabled
ratios are about 0.0028, 0.0031 and 0.0040; the delay dominates despite noisy
submillisecond A/A observations. The empty diagram is recognized before the
stage in both arms. These actual costs justify retaining the disabled default;
no reliable full-recognition speedup is claimed.

Runtime release \code{fbd3c9a8d}, the tests, exact source audit, paired source
and diagram timings are retained under \code{data/pachner-native-*}.
The driver is \code{causal\_research.native}, with \code{audit},
\code{source-benchmark} and \code{diagram-benchmark} modes.
Publication review checks 641 runtime pins, all completed timing outcomes,
all 124 source endpoint proofs and both strict descents, and the saved
source-bound diagram proof with move/search/transport producers disabled.
The maintained bounded primitive and optional source-bound disc adapter are
concrete integrations of report 79; the global structural closure hypothesis
and general quasi-polynomial recognition remain unproved.
